\section{Introduction}\label{sec:intro}

AlphaGo Zero~\citep{silver2017mastering} and AlphaZero~\citep{silver2018general} learned
superhuman Go from self-play alone. They did it with enormous compute: about 41 TPU-years
for AlphaZero's main Go run, and about 74 GPU-years for the ELF OpenGo
replication~\citep{tian2019elf}, by the estimates of \citet{wu2019accelerating}.
\citeauthor{wu2019accelerating}'s \katago{} paper, the reference point throughout this
report, closed much of that gap. It surpassed ELF's final model after 19~days on fewer
than 30 V100 GPUs, a reduction of about $50\times$. It got there through a set of mostly
general techniques: playout cap randomization, policy target pruning, global pooling, an
auxiliary opponent-policy target, and auxiliary ownership and score targets.

\gofer{} is a small Go engine written in Go. It uses Chinese area scoring, MCTS with a
PUCT selection rule, a GTP interface, and ONNX neural evaluation. Its README cites
\katago{} as its inspiration. Until this work, its learning loop (``v3'') was a shell
script. The script generated self-play JSONL on a CPU cloud instance, fine-tuned a
6-block, 64-channel residual network, and promoted a candidate after an arena of up to 200 games
against the champion. Two constraints shape this report:
\begin{enumerate}
  \item \textbf{No compute budget.} The owner cannot rent a GPU today. The pipeline therefore has
        to do something useful on free CPU resources, and it has to make a future rented
        GPU, most likely a cheap preemptible instance, as productive as possible from the
        first hour.
  \item \textbf{No strength results yet.} Without a real training run, we cannot honestly claim
        that the v4 pipeline produces a stronger engine than v3, let alone one comparable to
        \katago{}. Every claim below is limited to what we measured or can derive, and
        \cref{sec:limitations} lists what remains untested.
\end{enumerate}

\paragraph{What ``better'' means here.}
The v4 pipeline improves on its predecessor in four concrete ways, and it differs from the
\katago{} paper's recipe in several deliberate ways:
\begin{itemize}
  \item \textbf{Correct targets.} v3 stored ownership in absolute colour while every other
        target is from the side to move. Its \code{policy\_opp} field held the
        \emph{previous} move's policy rather than the opponent's reply. It split
        validation by row, so near-identical positions from one game landed on both sides
        of the split. All three are fixed, and tests cover each fix
        (\cref{sec:selfplay,sec:learner}).
  \item \textbf{Adoption of \katago{}'s general techniques at small scale.} We adopt
        playout-cap-aware training, global-pooling bias blocks, the auxiliary
        opponent-policy head, ownership and score heads, and the power-law replay window
        (\cref{sec:learner,sec:orchestration}). One change is deliberate:
        \katago{} \emph{discards} fast-search positions, while we keep them and mask only
        the policy loss. \Cref{sec:learner-masking} discusses the correlation risk this
        brings.
  \item \textbf{Cheaper iteration.} The data is about $40\times$ smaller, with no JSON
        parsing (\cref{sec:selfplay}), and self-play for the next cycle overlaps training.
  \item \textbf{Controlled promotion errors.} Gating uses a sequential probability ratio
        test~\citep{wald1945sequential}. It stops early on clearly worse or clearly better
        candidates, and it cuts the false-promotion rate at zero true gain from 8.0\% for
        v3 to 2.4\%, while promoting a $+50$ Elo candidate 90\% of the time against v3's
        67\%. The price is \emph{longer} gates near zero gain
        (\cref{sec:orchestration}), so this is an error-control gain, not a cost saving.
  \item \textbf{Eight defects, found by building the gate.} Five in the search: a stale
        pointer that stopped it leaving the root on $9\times9$; a first-play-urgency
        constant that collapsed the visit distribution---the policy training target---onto
        one move; forced playouts that outnumbered the search budget and flattened the same
        target; an inverted value sign that made more search strictly weaker; and a search
        that never saw a game end. Three in the measurement apparatus, which is why the
        others survived: correlated per-game seeds, a strength baseline that silently
        enhanced both players, and matches that stopped on the statistic under test. On the project's own
        reproducible baseline, the side searching three times harder went from 7 wins in
        200 games to 136, a sign flip from $-576$ to $+131$ Elo
        (\cref{sec:pipeline-found}). This is the strongest argument in this report
        for building the measurement apparatus before spending money on compute.
  \item \textbf{Robustness to cheap hardware.} Every stage can resume after a crash or
        preemption, the trainer can continue mid-run, and champions are immutable and
        never pruned (\cref{sec:orchestration,sec:infra}). A distillation path lets the
        project change network architecture without discarding the champion's knowledge
        (\cref{sec:learner-distill}). That cost was the reason an earlier ablation kept the
        6$\times$64 network.
\end{itemize}

\paragraph{Contributions.}
(i)~A complete, tested description of a small-scale \katago{}-style pipeline for
$9\times9$ Go, from Go self-play to a gated, published ONNX champion. (ii)~An explicit
technique-by-technique comparison with \citet{wu2019accelerating}
(\cref{tab:overview-compare}). (iii)~Measurements of what can be measured without a
training run: data size, load time, parameter and compute budgets, CPU inference latency,
export parity, exact gate operating characteristics, and end-to-end cycle timings
(\cref{sec:eval}), together with the engine defects that apparatus exposed
(\cref{sec:pipeline-found}).
(iv)~An ablation protocol, modelled on \katago{}'s, for the experiments that would settle
the open questions (\cref{sec:conclusion}).

\paragraph{Organisation.}
\Cref{sec:background} reviews AlphaZero-style training and the relevant parts of
\katago{}. \Cref{sec:overview} gives the system overview and comparison tables.
\Cref{sec:selfplay,sec:learner,sec:orchestration,sec:infra} describe self-play data, the
learner, orchestration, and infrastructure in turn. \Cref{sec:eval} reports measurements,
\cref{sec:limitations} lists limitations, and \cref{sec:conclusion} concludes.
